% Stage25 manuscript integration

\subsection{Prevalence and Operational Stress Results}

Stage25 translated the already-frozen Stage22 and Stage24 operating
characteristics into deployment-stress quantities without model
refitting, new inference, target reopening, threshold re-selection, or
calibration. Twenty-four frozen operating points were evaluated across
six preregistered attack prevalences (10\%, 3\%, 1\%, 0.3\%, 0.1\%,
and 0.01\%), yielding 144 deterministic prior-shift projections.

At 0.1\% prevalence, the Stage22 random-natural STANDARD operating point
retained PPV 0.965572, with
34.6 false alerts/day and
969.4 true alerts/day per one million
benign flows. The projected processing requirement was
33.5
analyst-hours/day. The SECURITY operating point increased frozen recall
to 0.984701, but its FPR of
0.006817 produced
6817.1 false alerts/day, reducing PPV
to 0.126325.

The chronological STANDARD operating point projected PPV
0.000551068 and only
0.0322 true detections/day at 0.1\%
prevalence. For IDS2018$\rightarrow$CICIDS2017 STANDARD transfer,
projected PPV ranged from 0.039233 to
0.060313, whereas the reverse
CICIDS2017$\rightarrow$IDS2018 STANDARD transfer yielded only
0.000257610--0.000287993.

Under the frozen relative-cost ratio $C_{FP}:C_{FN}=1:100$, 15/24
operating points had lower relative cost than the simplified ignore
reference at 0.1\% prevalence, compared with only 3/24 at 0.01\%.
All seven preregistered sanity tests passed.

\subsection{Discussion}

Stage25 separates prior-probability shift from temporal and domain shift.
The former changes PPV even when TPR and FPR are held fixed, whereas the
latter can change the operating characteristics themselves. The results
show that low FPR is a critical requirement for preserving alert
precision at rare attack prevalences, and that high benchmark recall can
still produce an operationally excessive false-alert stream.

SOC capacity and detection usefulness must also be interpreted
separately. An operating point can exceed a small analyst team's capacity
despite high PPV because true alerts require service, while another can
fit within capacity simply because detection has collapsed. The
bidirectional Stage24 asymmetry remains visible under this operational
translation.

The analytic break-even calculations identify exact prevalence
boundaries under the frozen assumptions and avoid post-hoc graphical
selection. They are conditional decision-analysis tools rather than
claims of universal deployment readiness.

\subsection{Limitations and Threats to Validity}

Stage25 assumes prior-probability shift within each frozen operating
point and is not empirical production validation. Real environments may
change the class-conditional feature distributions and therefore TPR and
FPR. The fixed one-million-benign-flow volume, two-minute alert service
time, and one/three/ten analyst-day capacity tiers are reference
scenarios rather than universal SOC constants.

The $C_{FP}=1$, $C_{FN}=100$ cost model uses dimensionless relative
operational cost units and must not be interpreted as financial loss.
The projections are deterministic conditional transformations of frozen
empirical estimates; TPR/FPR estimation uncertainty is not propagated.
The benchmark datasets also differ from contemporary live enterprise
telemetry. Finally, Stage24 GROUNDED\_S4 cells remained non-evaluable
because exact durable physical membership was unavailable, and Stage25
does not introduce a heuristic replacement.

\subsection{Contributions}

Stage25 provides a validation-safe operational translation layer that:
(i) projects frozen operating points across preregistered attack
prevalences using exact Bayesian relations; (ii) converts operating
characteristics into projected alert volume and SOC workload;
(iii) derives analytic PPV, required-FPR, and relative-cost break-even
boundaries; (iv) connects random, chronological, and bidirectional
cross-dataset results without averaging incompatible directions; and
(v) preserves preregistered assumptions, figures, sanity tests, and
artifact-level reproducibility through the final scientific seal.
